% phu-macros.tex: the one list of math macros.
%
% Read by phumacros.sty (LaTeX: papers, reports, Quarto PDF) and by the
% Quarto `phu` filter, which hands them to MathJax for HTML. So the same
% \R or \argmax works in every output.
%
% Keep it to one definition per line, using only
%   \newcommand  \renewcommand  \providecommand  \DeclareMathOperator(*)
% so the filter can read it. Text-mode helpers that MathJax can't use
% belong in phunotes.sty instead.

% --- code output (matrices written by phu.py / phu.R) ----------------
\providecommand{\phuneg}[1]{\textcolor[HTML]{B4442B}{#1}}

% --- number sets ------------------------------------------------------
\providecommand{\R}{\mathbb{R}}
\providecommand{\C}{\mathbb{C}}
\providecommand{\N}{\mathbb{N}}
\providecommand{\Z}{\mathbb{Z}}
\providecommand{\Q}{\mathbb{Q}}

% --- delimiters -------------------------------------------------------
\providecommand{\abs}[1]{\left\lvert #1 \right\rvert}
\providecommand{\norm}[1]{\left\lVert #1 \right\rVert}
\providecommand{\mean}[1]{\left\langle #1 \right\rangle}
\providecommand{\sprod}[2]{\left\langle \boldsymbol{#1},\boldsymbol{#2}\right\rangle}
\providecommand{\set}[1]{\left\{ #1 \right\}}

% --- vectors and derivatives (the bits of `physics` actually used) ----
\providecommand{\vb}[1]{\boldsymbol{#1}}
\renewcommand{\vec}[1]{\boldsymbol{#1}}
\providecommand{\dd}{\mathrm{d}}
\providecommand{\dv}[2]{\frac{\mathrm{d}#1}{\mathrm{d}#2}}
\providecommand{\pdv}[2]{\frac{\partial #1}{\partial #2}}

% --- probability ------------------------------------------------------
\providecommand{\E}{\mathbb{E}}
\providecommand{\Prob}{\mathbb{P}}
\providecommand{\Var}{\operatorname{Var}}
\providecommand{\Cov}{\operatorname{Cov}}
\providecommand{\indep}{\perp\!\!\!\perp}

% --- operators --------------------------------------------------------
\DeclareMathOperator*{\argmax}{arg\,max}
\DeclareMathOperator*{\argmin}{arg\,min}
\DeclareMathOperator{\Tr}{Tr}
\DeclareMathOperator{\sgn}{sgn}

% --- colour emphasis --------------------------------------------------
\providecommand{\red}[1]{\textcolor{red}{#1}}
\providecommand{\blue}[1]{\textcolor{blue}{#1}}
\providecommand{\green}[1]{\textcolor{green}{#1}}
\providecommand{\orange}[1]{\textcolor{orange}{#1}}
